% GENERATED FILE. Edit canonical lessons, then run npm run generate:books.
% canonical-source-hash: fnv1a64:db240ee7fdd589b7
% canonical-lessons: FR-C335-debat, FR-C335-debattre, FR-C335-these, FR-C335-argumentation, FR-C335-persuader

\chapter{Le débat}
\label{ch:fr-le-debat}
\hlchaptermodality{335}

\begin{chapteropening}
\textbf{By the end of this chapter:} \emph{I can take part in a debate: debate a question, state a thesis, build a line of argument and try to persuade someone.}

Complete the last lesson of chapter 335: I can take part in a debate: debate a question, state a thesis, build a line of argument and try to persuade someone.
\end{chapteropening}

\section[le débat]{le débat — the debate}

\label{lesson:FR-C335-debat}



\begin{quote}
Before the new one: say the French for to state firmly, then the French for convincing.
\end{quote}

\subsection*{You'll want to know: le débat}
\textbf{le débat} — \textquotedblleft{}the debate\textquotedblright{}.

\begin{quote}
Ce soir, il y a un débat à la télévision sur l'école.
\end{quote}

\begin{quote}
\emph{Tonight there is a debate on television about schools.}
\end{quote}

\emph{Un débat animé} — a lively debate.

\begin{cousinweb}
\emph{débat} is formed from the verb \emph{débattre}, built on \emph{battre} \textquotedblleft{}to beat\textquotedblright{}, which is inherited from Latin \emph{battuere} \textquotedblleft{}to beat, to strike\textquotedblright{} (EtymDB; Lewis \& Short). Ideas are knocked back and forth.
\end{cousinweb}

\begin{grammarlens}[title={What you've built}]
One more word for telling, explaining and describing at B1.
\end{grammarlens}

\subsection*{Your turn}
\begin{itemize}
\raggedright
  \item \emph{Say it:} \emph{le débat}
  \item \emph{Say it:} \emph{le débat}, once more
  \item \emph{Say it:} a sentence of your own with \emph{le débat}
  \item \emph{Recall:} say the French for to state firmly, then the French for convincing, then say \emph{le débat} again
\end{itemize}

\begin{tcolorbox}[breakable,colback=teal!4,colframe=teal!35!black,title={Before you move on}]
What does le débat mean? (\textquotedblleft{}The debate\textquotedblright{}.) Say it once more.
\end{tcolorbox}
\section[débattre]{débattre — to debate, to discuss}

\label{lesson:FR-C335-debattre}



\begin{quote}
Before the new one: say the French for convincing, then the French for the debate.
\end{quote}

\subsection*{You'll want to know: débattre}
\textbf{débattre} — \textquotedblleft{}to debate, to discuss\textquotedblright{}.

\begin{quote}
Nous avons débattu de ce problème pendant deux heures.
\end{quote}

\begin{quote}
\emph{We debated this problem for two hours.}
\end{quote}

\emph{Débattre de quelque chose}. Conjugated like \emph{battre}: \emph{je débats, nous débattons}.

\begin{cousinweb}
\emph{débattre} is inherited from Old French \emph{debatre} \textquotedblleft{}to fight, to contend\textquotedblright{} (EtymDB), built from \emph{dé-} and \emph{battre}, which comes from Latin \emph{battuere} \textquotedblleft{}to beat\textquotedblright{} (Lewis \& Short).
\end{cousinweb}

\begin{grammarlens}[title={What you've built}]
One more word for telling, explaining and describing at B1.
\end{grammarlens}

\subsection*{Your turn}
\begin{itemize}
\raggedright
  \item \emph{Say it:} \emph{débattre}
  \item \emph{Say it:} \emph{débattre}, once more
  \item \emph{Say it:} a sentence of your own with \emph{débattre}
  \item \emph{Recall:} say the French for convincing, then the French for the debate, then say \emph{débattre} again
\end{itemize}

\begin{tcolorbox}[breakable,colback=teal!4,colframe=teal!35!black,title={Before you move on}]
What does débattre mean? (\textquotedblleft{}To debate, to discuss\textquotedblright{}.) Say it once more.
\end{tcolorbox}
\section[la thèse]{la thèse — the thesis, the position argued for}

\label{lesson:FR-C335-these}



\begin{quote}
Before the new one: say the French for the debate, then the French for to debate.
\end{quote}

\subsection*{You'll want to know: la thèse}
\textbf{la thèse} — \textquotedblleft{}the thesis, the position argued for\textquotedblright{}.

\begin{quote}
Sa thèse est simple : nous travaillons trop.
\end{quote}

\begin{quote}
\emph{His thesis is simple: we work too much.}
\end{quote}

At university, \emph{une thèse} is a doctoral thesis.

\begin{cousinweb}
\emph{thèse} is a learned borrowing from Latin \emph{thesis} \textquotedblleft{}a proposition\textquotedblright{} (Lewis \& Short), which came from Greek \emph{thésis} \textquotedblleft{}a setting down, a statement\textquotedblright{} (EtymDB).
\end{cousinweb}

\begin{grammarlens}[title={What you've built}]
One more word for telling, explaining and describing at B1.
\end{grammarlens}

\subsection*{Your turn}
\begin{itemize}
\raggedright
  \item \emph{Say it:} \emph{la thèse}
  \item \emph{Say it:} \emph{la thèse}, once more
  \item \emph{Say it:} a sentence of your own with \emph{la thèse}
  \item \emph{Recall:} say the French for the debate, then the French for to debate, then say \emph{la thèse} again
\end{itemize}

\begin{tcolorbox}[breakable,colback=teal!4,colframe=teal!35!black,title={Before you move on}]
What does la thèse mean? (\textquotedblleft{}The thesis, the position argued for\textquotedblright{}.) Say it once more.
\end{tcolorbox}
\section[l'argumentation]{l'argumentation — the line of argument}

\label{lesson:FR-C335-argumentation}



\begin{quote}
Before the new one: say the French for to debate, then the French for the thesis.
\end{quote}

\subsection*{You'll want to know: l'argumentation}
\textbf{l'argumentation} — \textquotedblleft{}the line of argument\textquotedblright{}.

\begin{quote}
Ton argumentation est claire, mais je ne suis pas d'accord.
\end{quote}

\begin{quote}
\emph{Your line of argument is clear, but I do not agree.}
\end{quote}

The whole chain of reasons; one single reason is \emph{un argument}. Feminine.

\begin{cousinweb}
\emph{argumentation} is a learned borrowing from Latin \emph{argūmentātiō} \textquotedblleft{}the bringing of proof\textquotedblright{} (EtymDB; Lewis \& Short), from \emph{argūmentum} \textquotedblleft{}proof\textquotedblright{}, built on \emph{arguere} \textquotedblleft{}to make clear\textquotedblright{}.
\end{cousinweb}

\begin{grammarlens}[title={What you've built}]
One more word for telling, explaining and describing at B1.
\end{grammarlens}

\subsection*{Your turn}
\begin{itemize}
\raggedright
  \item \emph{Say it:} \emph{l'argumentation}
  \item \emph{Say it:} \emph{l'argumentation}, once more
  \item \emph{Say it:} a sentence of your own with \emph{l'argumentation}
  \item \emph{Recall:} say the French for to debate, then the French for the thesis, then say \emph{l'argumentation} again
\end{itemize}

\begin{tcolorbox}[breakable,colback=teal!4,colframe=teal!35!black,title={Before you move on}]
What does l'argumentation mean? (\textquotedblleft{}The line of argument\textquotedblright{}.) Say it once more.
\end{tcolorbox}
\section[persuader]{persuader — to persuade}

\label{lesson:FR-C335-persuader}



\begin{quote}
Before the new one: say the French for the thesis, then the French for the line of argument.
\end{quote}

\subsection*{You'll want to know: persuader}
\textbf{persuader} — \textquotedblleft{}to persuade\textquotedblright{}.

\begin{quote}
J'ai essayé de persuader mon père d'acheter une nouvelle voiture.
\end{quote}

\begin{quote}
\emph{I tried to persuade my father to buy a new car.}
\end{quote}

\emph{Persuader quelqu'un de faire quelque chose}. Close to \emph{convaincre}.

\begin{cousinweb}
\emph{persuader} comes from Latin \emph{persuādēre} \textquotedblleft{}to bring over by talking, to convince\textquotedblright{} (Lewis \& Short), through Middle French (EtymDB). It is built from \emph{per-} and \emph{suādēre} \textquotedblleft{}to advise\textquotedblright{}.
\end{cousinweb}

\begin{grammarlens}[title={What you've built}]
One more word for telling, explaining and describing at B1.
\end{grammarlens}

\subsection*{Your turn}
\begin{itemize}
\raggedright
  \item \emph{Say it:} \emph{persuader}
  \item \emph{Say it:} \emph{persuader}, once more
  \item \emph{Say it:} a sentence of your own with \emph{persuader}
  \item \emph{Recall:} say the French for the thesis, then the French for the line of argument, then say \emph{persuader} again
\end{itemize}

\begin{tcolorbox}[breakable,colback=teal!4,colframe=teal!35!black,title={Before you move on}]
What does persuader mean? (\textquotedblleft{}To persuade\textquotedblright{}.) Say it once more.
\end{tcolorbox}
